\documentclass[12pt,a4paper]{article}

\usepackage{amsmath,amssymb,amsthm}
\usepackage{hyperref}
\usepackage{natbib}
\usepackage{geometry}
\usepackage{booktabs}
\usepackage{xcolor}
\usepackage{microtype}
\usepackage{float}
\usepackage{mdframed}
\usepackage{array}
\usepackage{colortbl}

\geometry{margin=1in}

\definecolor{mcmcrow}{RGB}{220,230,250}
\definecolor{newresult}{RGB}{200,230,200}

\hypersetup{
  colorlinks=true,
  linkcolor=blue,
  citecolor=blue,
  urlcolor=blue,
  pdftitle={IAM's Law: The Thermodynamic Cost of Classical Existence}
}

\newtheorem{law}{Law}
\newtheorem{proposition}{Proposition}
\newtheorem{corollary}{Corollary}
\newtheorem{definition}{Definition}

% ─────────────────────────────────────────────────────────────────────────────
\title{\textbf{IAM's Law}\\[0.5em]
\large The Thermodynamic Cost of Classical Existence}

\author{Heath W.\ Mahaffey\\
\small Independent Researcher, Entiat, Washington 98822, USA\\
\small \href{mailto:hmahaffeyges@gmail.com}{hmahaffeyges@gmail.com}\\
\small \url{https://github.com/hmahaffeyges/IAM-Validation}}

\date{March 2026\\
\small Zenodo DOI:
\href{https://doi.org/10.5281/zenodo.18702042}{10.5281/zenodo.18702042}}
% ─────────────────────────────────────────────────────────────────────────────

\begin{document}
\maketitle

% ─────────────────────────────────────────────────────────────────────────────
\begin{abstract}
We present and derive a candidate thermodynamic law governing the
cost of classical existence. Every irreversible transition from
quantum potential to classical actuality --- every decoherence event
on a timelike worldline --- pays a thermodynamic cost of $k_B T_H
\ln 2$ per bit to the nearest encoding surface. This cost is
physically real, irreversible, and cosmologically consequential.
We call this IAM's Law.

The law has two expressions at two scales. Locally, it provides the
first complete physical interpretation of the virial theorem,
established by Clausius in 1870~\citep{Clausius1870}: the kinetic
energy $\langle K \rangle$ of any $1/r$ bound system is the Landauer
cost of the system's transition from unbound potential to bound
actuality. The virial theorem $2\langle K \rangle + \langle V
\rangle = 0$ is the mechanical face of this identity. The
thermodynamic completion $\langle K \rangle = Q = T\,\Delta S =
E_\mathrm{Landauer}$ is its other face. The $1/2$ partition is the
unique ratio at which mechanical and thermodynamic equilibrium are
simultaneously satisfied for a $1/r$ potential.

Globally, the accumulated Landauer cost of 13.8 billion years of
decoherence events constitutes a missing term in the
Jacobson--Cai-Kim entropy functional. Adding it produces the
activation function $\mathcal{E}(a) = \exp(1 - 1/a)$, a
sector-dependent gravitational coupling $\mu(a) < 1$ for matter
and $\Sigma(a) = 1$ for photons, and a coupling constant $\beta_m
= \Omega_m/2$ derived from the virial theorem with zero free
parameters. This cosmological model is one derived implementation
of the broader law; the law itself operates locally through the
virial theorem and globally through the horizon entropy completion.

The law is applied across 37 orders of magnitude with one equation
and zero free parameters. The cosmological prediction $\beta_m =
\Omega_m/2 = 0.1575$ is consistent with 18 independent converged
MCMC chains at $0.2\sigma$~\citep{PlanckVI2020}: 17 testing the
IAM growth equation against the full Planck~2018 likelihood, and
one dedicated baryon asymmetry chain returning $\eta =
(6.115 \pm 0.037)\times10^{-10}$ with the BBN prior removed.
\end{abstract}

\noindent\textbf{Keywords:} IAM's Law; thermodynamic cost of actuality;
virial theorem; Landauer's principle; gravitational decoherence;
Jacobson--Cai-Kim entropy; horizon thermodynamics; dark energy;
$\sigma_8$ tension; $H_0$ tension

% ─────────────────────────────────────────────────────────────────────────────
\section{The Law in Plain Language}
\label{sec:plain}
% ─────────────────────────────────────────────────────────────────────────────

Before the mathematics, a plain statement of what this paper claims.

The universe contains two kinds of things: things that exist
classically --- with definite positions, definite states, definite
identities --- and things that exist only as quantum potential ---
superpositions of possibilities that have not yet resolved into definite
outcomes. The transition from the second kind to the first kind is
called decoherence. It happens constantly, everywhere: every atom that
binds, every molecule that forms, every star that collapses, every
galaxy that virializes.

Here is the claim: that transition is not free. Every time something
crosses from quantum potential to classical actuality, the universe
charges a thermodynamic price. That price is paid irreversibly to the
boundary of the observable universe --- the cosmic horizon --- at a rate
of $k_B T_H \ln 2$ per bit of classical information produced. Once
paid, the cost cannot be recovered. Classical existence is continuously
purchased, one decoherence event at a time.

This is not a metaphor. It is a physical claim with a derivation, a
mechanism, and measured consequences.

\textbf{Locally,} the virial theorem $2\langle K \rangle + \langle V
\rangle = 0$ has governed the energy of bound physical systems since
Clausius stated it in 1870~\citep{Clausius1870}. For any system bound by
a $1/r$ potential --- a hydrogen atom, a solar system, a galaxy --- the
kinetic energy at equilibrium is exactly half the magnitude of the
potential energy. This is one of the most broadly verified results in
physics, confirmed from the electron orbital to the cosmic web. The
virial theorem has historically been stated as a mechanical result,
without a thermodynamic interpretation of $\langle K \rangle$ or an
account of where the other half of $|\langle V \rangle|$ goes when a
system virializes.

IAM's Law answers it. The other half of $|\langle V \rangle|$ is the
heat $Q$ released irreversibly at the virialization boundary --- the
thermodynamic cost of the system's transition from unbound potential to
bound actuality. By Landauer's principle~\citep{Landauer1961}, this cost
equals the kinetic energy retained in the bound state:
$Q = \langle K \rangle = E_\mathrm{Landauer}$. The virial theorem is the
mechanical face of this identity. The $1/2$ partition is not arbitrary.
It is the unique ratio at which mechanical and thermodynamic equilibrium
are simultaneously satisfied for a $1/r$ potential. One coin, two sides.

\textbf{Globally,} Jacobson showed in 1995 that the Einstein field
equations follow from thermodynamics applied to local causal
horizons~\citep{Jacobson1995}. Cai and Kim extended this to the FRW
apparent horizon in 2005, recovering the Friedmann
equations~\citep{CaiKim2005}. Both derivations assume the entropy
functional of the horizon contains only geometric terms. IAM's Law says
that assumption is incomplete: the Landauer cost of every irreversible
decoherence event in the bulk must also be encoded on the horizon. This
is the term that Einstein's GR lacked. It is what Jacobson was one
entropy term short of. Adding it completes the derivation that Jacobson
began.

The virial theorem then serves as the balance sheet at the cosmological
scale: it tells you exactly how much energy crosses the actualization
boundary at each virialization event, giving $\beta_m = \Omega_m/2$ with
no free parameters. Seventeen converged MCMC chains against the full
Planck~2018 likelihood are consistent with this prediction at $0.2\sigma$.

The law is the same at both scales. The cost is the same. The balance
sheet is the virial theorem. The only difference is the scale at which
the ledger is read.

The virial theorem, Landauer's principle, and the second law of
thermodynamics are three faces of the same underlying asymmetry: the
cost of becoming is paid in one direction only. Potential does not
reclaim what actuality has spent. This is why $\mathcal{E}(a)$ is
monotonically increasing and never reaches $e$ --- the universe
approaches but never completes its own actualization. IAM's Law is a
statement about the cost of that process at every scale where it occurs.

Landauer's principle is the special case of IAM's Law applied to
computational information erasure at a fixed laboratory temperature.
IAM's Law is the governing principle of which Landauer's principle is
one expression --- applied here to any irreversible actualization event
at the local horizon temperature, regardless of the physical system or
scale. The virial theorem is a second expression of the same law,
governing the equilibrium partition when the actualization process has
settled into a bound state. The second law is a third expression,
governing the direction of the accumulated cost. One asymmetry. Three
faces.

The temperature distinction establishes the hierarchy precisely.
In Landauer's principle, the temperature $T$ of the thermal reservoir
is an external free parameter: the principle states the minimum cost
is $k_B T \ln 2$ per bit, but $T$ must be supplied from outside the
system. The principle is silent on where $T$ comes from. In IAM's Law,
the temperature $T_H$ is not free. It is derived from the geometry of
the nearest encoding surface --- the Gibbons-Hawking temperature
$T_H = \hbar H / (2\pi k_B)$ at the cosmic horizon, or the Hawking
temperature $T_H = \hbar c^3 / (8\pi G M k_B)$ at a black hole
horizon. The temperature is determined by the physics of the boundary
itself, not supplied externally.

This is the basis for the hierarchy claim. A principle that derives
its temperature parameter from the geometry of the encoding surface
encompasses a principle that requires that parameter to be supplied
as input. Landauer's principle is recovered as the special case where
the nearest encoding surface is the local thermal environment at
laboratory temperature --- the limit in which $T_H \to T_\mathrm{lab}$.

The virial derivation in Section~\ref{sec:local} makes this concrete.
The temperature cancels exactly in Step~3 (Eq.~\eqref{eq:T_cancels}),
producing a result that holds \textit{independent of temperature and
independent of scale}. Landauer's principle, by construction, cannot
achieve this: it requires $T$ to be known. IAM's Law does not. The
temperature-independence of the virial identity is the derivational
proof that IAM's Law is the more fundamental statement.

% ─────────────────────────────────────────────────────────────────────────────
\section{Statement of the Law}
\label{sec:statement}
% ─────────────────────────────────────────────────────────────────────────────

\begin{mdframed}
\begin{law}[IAM's Law]
\label{law:iam}
Every irreversible transition from quantum potential to classical
actuality --- every decoherence event on a timelike worldline ---
pays a thermodynamic cost of $k_B T_H \ln 2$ per bit of classical
information produced, encoded permanently on the nearest available
encoding surface at the local horizon temperature $T_H$:
\begin{equation}
\Delta E_\mathrm{act} = k_B T_H \ln 2 \quad \text{per bit.}
\label{eq:iams_law}
\end{equation}
The transition is irreversible. The cost cannot be recovered.
Potential moves toward actual; never the reverse.
\end{law}
\end{mdframed}

\medskip

Three consequences follow immediately:

\begin{corollary}[The thermodynamic arrow of time]
Each payment of the actualization cost is irreversible. The
accumulated ledger of all such payments, tracked by
$\mathcal{E}(a)$, is monotonically increasing. The thermodynamic
arrow of time is a structural property of IAM's Law, derived
quantitatively in Section~\ref{sec:arrow}.
\end{corollary}

% ─────────────────────────────────────────────────────────────────────────────
\section{The Hierarchy}
\label{sec:hierarchy}
% ─────────────────────────────────────────────────────────────────────────────

IAM's Law generates a complete hierarchy at two scales.

\medskip
\noindent\textbf{IAM's Law} [Section~\ref{sec:statement}] ---
the root principle: every crossing pays $k_B T_H \ln 2$ per bit.

\medskip
\noindent\textbf{Local expression} [Section~\ref{sec:local}]:
IAM's Law applied to $1/r$ bound systems.
\begin{itemize}
\item \textit{The virial theorem derived}: the equilibrium condition
  $2\langle K \rangle + \langle V \rangle = 0$ follows from Euler's
  theorem for homogeneous functions of degree $-1$, not assumed.
\item \textit{Thermodynamic completion}: $\langle K \rangle = Q =
  T\,\Delta S = E_\mathrm{Landauer}$. The physical meaning of $K$,
  identified for the first time.
\item \textit{The $1/2$ partition}: the unique ratio at which
  mechanical and thermodynamic equilibrium are simultaneously
  satisfied for $1/r$. Verified across 37 orders of magnitude.
\end{itemize}

\medskip
\noindent\textbf{Global expression} [Section~\ref{sec:global}]:
IAM's Law integrated across cosmological time.
\begin{itemize}
\item \textit{The missing entropy term}: accumulated Landauer cost
  of 13.8 billion years of decoherence events, encoded on the
  cosmic horizon --- the term Jacobson and Cai-Kim lacked.
\item \textit{The activation function}: $\mathcal{E}(a) =
  \exp(1-1/a)$, derived from horizon thermodynamics.
\item \textit{The virial theorem as balance sheet}: $\beta_m =
  \Omega_m/2$, connecting the local $1/2$ partition to the
  cosmological coupling constant. Zero free parameters.
\item \textit{The dual-sector split}: $\mu(a) < 1$ for matter
  on timelike worldlines; $\Sigma(a) = 1$ for photons on null
  geodesics. Geometric necessity, not assumption.
\end{itemize}

\medskip
\noindent\textbf{The bridge} [Section~\ref{sec:bridge}]: the critical
exponent $n = 5/2$, derived independently from both directions.
Their convergence is a nontrivial consistency check that the same
law operates at both scales.

% ─────────────────────────────────────────────────────────────────────────────
\section{The Quantum Boundary of Reality}
\label{sec:quantum}
% ─────────────────────────────────────────────────────────────────────────────

IAM's Law operates at the boundary between quantum and classical. To
state it precisely, we must state what that boundary is.

\subsection{The Decoherence Event}

Consider a quantum system $\mathcal{S}$ coupled to an environment
$\mathcal{E}$. An initial product state evolves under the total
Hamiltonian into an entangled state~\citep{Zurek2003}:
\begin{equation}
|\psi_S\rangle \otimes |E_0\rangle \;\longrightarrow\;
\sum_i c_i\,|s_i\rangle \otimes |E_i\rangle,
\label{eq:decoherence}
\end{equation}
where $\{|s_i\rangle\}$ are the pointer states selected by the
system-environment interaction. When the environmental states become
approximately orthogonal ($\langle E_i | E_j \rangle \approx
\delta_{ij}$), the reduced density matrix of $\mathcal{S}$ becomes
diagonal:
\begin{equation}
\rho_S = \mathrm{Tr}_E\,|\Psi\rangle\langle\Psi|
\;\longrightarrow\; \sum_i |c_i|^2\,|s_i\rangle\langle s_i|.
\label{eq:rho_diagonal}
\end{equation}
Coherence is lost. The system has crossed from quantum superposition
to classical actuality. This crossing is irreversible: recovering the
off-diagonal elements requires accessing all environmental correlations,
which IAM's Law --- of which Landauer's principle~\citep{Landauer1961}
is the special case at fixed laboratory temperature ---
renders thermodynamically prohibited at any finite temperature.

\subsection{Information Produced at the Boundary}

Each decoherence event produces classical information: the specification
of which pointer state $|s_i\rangle$ was selected. For a system
decohering into $N$ distinguishable pointer states:
\begin{equation}
I = \log_2 N \quad \text{bits.}
\label{eq:info_produced}
\end{equation}
Landauer's principle establishes that producing (equivalently,
irreversibly encoding) one bit of classical information requires a
minimum energy dissipation of:
\begin{equation}
\Delta E_\mathrm{Landauer} = k_B T \ln 2
\label{eq:landauer_basic}
\end{equation}
at temperature $T$. This bound has been experimentally
verified~\citep{Berut2012,Jun2014}. The irreversibility is not
practical but thermodynamic: undoing the decoherence would require
erasing the produced information, dissipating energy into the
environment and increasing its entropy. The second law forbids
the reverse.

\subsection{The Nearest Encoding Surface}

IAM's Law specifies that the cost is paid to the \textit{nearest
available encoding surface} at its local temperature $T_H$. In the
local regime, for gravitationally collapsing systems, the nearest
encoding surface is the Bekenstein-Hawking horizon of the forming
black hole, at Hawking temperature $T_\mathrm{BH} = \hbar c^3 /
(8\pi G M k_B)$. In the cosmological regime, the encoding surface
is the Gibbons-Hawking cosmic horizon at temperature $T_H = \hbar H /
(2\pi k_B)$~\citep{Gibbons1977}. Both are thermal surfaces. The same
law governs both. The scale determines which surface receives the cost:
concentrated local collapses write first to the nearest black hole
horizon; distributed large-scale structure formation writes to the
growing cosmic horizon as it expands and cools. The transition between
these two regimes is tracked quantitatively by the activation function
$\mathcal{E}(a)$ derived in Section~\ref{sec:global}.

\subsection{The Epoch-Dependent Cost}

At the cosmological horizon, the cost per bit is epoch-dependent:
\begin{equation}
\Delta E_\mathrm{act}(a) = k_B T_H(a) \ln 2
= \frac{\hbar H(a)}{2\pi}\,\ln 2.
\label{eq:epoch_cost}
\end{equation}
When $H$ is large (early universe), bits are expensive. When $H$ is
small (late universe), bits are cheap. The total cost paid over cosmic
history depends on both the rate of information production and the
varying cost per bit --- this is what gives rise to the activation
function $\mathcal{E}(a)$, derived in Section~\ref{sec:global}.

% ─────────────────────────────────────────────────────────────────────────────
\section{Local Expression: The Complete Virial Theorem}
\label{sec:local}
% ─────────────────────────────────────────────────────────────────────────────

\subsection{Setup}

Consider a system of total mass $M$ transitioning from an unbound state
to a stable bound equilibrium under a $1/r$ gravitational potential.
Natural reference state: unbound at rest at infinity, total energy
$E_i = 0$. The bound equilibrium state has total energy:
\begin{equation}
E_f = \langle K \rangle + \langle V \rangle.
\label{eq:Ef}
\end{equation}
No assumption about the ratio $\langle K \rangle / |\langle V \rangle|$
is made at this stage.

\subsection{Step 1 --- First Law: Heat Released at the Boundary}

The transition from unbound to bound is irreversible. By the first law:
\begin{equation}
Q = E_i - E_f = -E_f = -(\langle K \rangle + \langle V \rangle).
\label{eq:Q_firstlaw}
\end{equation}
This is exact. $Q > 0$ requires $E_f < 0$, satisfied for any stable
bound state.

\subsection{Step 2 --- Second Law: Entropy Produced}

The transition is irreversible. By the second law:
\begin{equation}
\Delta S = \frac{Q}{T_\mathrm{vir}},
\label{eq:entropy_step2}
\end{equation}
where $T_\mathrm{vir} = m\sigma^2 / 3k_B$ is the virial temperature
of the bound system, with $m$ the characteristic particle mass and
$\sigma$ the velocity dispersion.

\subsection{Step 3 --- IAM's Law Applied: The Actualization Cost (Landauer's Principle as Special Case)}

By IAM's Law, the minimum thermodynamic cost of this irreversible
transition is the cost of the classical information produced at the
actualization boundary. Landauer's principle~\citep{Landauer1961}
is the special case of IAM's Law at fixed laboratory temperature;
here, IAM's Law applies at the local horizon temperature $T_H$, which
cancels exactly in the derivation below. The actualization cost is:
\begin{equation}
E_\mathrm{Landauer} = T_\mathrm{vir}\,\Delta S.
\label{eq:landauer_step3}
\end{equation}
Substituting Eq.~\eqref{eq:entropy_step2}:
\begin{equation}
E_\mathrm{Landauer}
= T_\mathrm{vir} \times \frac{Q}{T_\mathrm{vir}} = Q.
\label{eq:T_cancels}
\end{equation}
The virial temperature cancels exactly. The actualization cost equals
the heat released at the virialization boundary
\textit{independent of temperature and independent of scale}. This
is the scale-invariance of IAM's Law: the same thermodynamic identity
holds from hydrogen atoms to galaxy clusters because the temperature
cancels at every scale.

\subsection{Step 4 --- Euler's Theorem: The Virial Theorem Derived}

Steps 1--3 establish:
\begin{equation}
E_\mathrm{Landauer} = Q = -E_f
= -(\langle K \rangle + \langle V \rangle).
\label{eq:combined_local}
\end{equation}
To complete the identification of $\langle K \rangle$, we invoke the
equilibrium condition of the $1/r$ potential. For any potential
$V \propto r^n$, Euler's theorem for homogeneous functions of degree
$n$ requires at stable equilibrium $2\langle K \rangle + n\langle V
\rangle = 0$. For $1/r$ potentials ($n = -1$):
\begin{equation}
2\langle K \rangle + \langle V \rangle = 0
\quad\Longrightarrow\quad
-E_f = -(\langle K \rangle + \langle V \rangle)
= \langle K \rangle.
\label{eq:euler_virial}
\end{equation}
This is the virial theorem --- not assumed, but derived from Euler's
theorem as the mathematical condition that defines stable equilibrium
under a $1/r$ potential. Any system not satisfying
Eq.~\eqref{eq:euler_virial} is either not yet at equilibrium or
is unstable.

\subsection{The Thermodynamic Completion}

Substituting Eq.~\eqref{eq:euler_virial} into
Eq.~\eqref{eq:combined_local}:
\begin{equation}
\boxed{E_\mathrm{Landauer} = Q = T_\mathrm{vir}\,\Delta S
= -E_f = \langle K \rangle
= \tfrac{1}{2}|\langle V \rangle|.}
\label{eq:identity}
\end{equation}
The chain is exact at every step. The virial theorem
$2\langle K \rangle + \langle V \rangle = 0$ is the
\textit{mechanical face}: the equilibrium ratio of the $1/r$
potential, known since Clausius~\citep{Clausius1870}. The
thermodynamic completion $\langle K \rangle = Q = T\,\Delta S
= E_\mathrm{Landauer}$ is the \textit{thermodynamic face}: the
physical identity of what $K$ is within the IAM framework.
These are not two separate results. They are the two faces of the same
partition of $|\langle V \rangle|$ at the actualization boundary.

\subsection{The Physical Meaning of K}

$\langle K \rangle$ is not merely the average kinetic energy of bound
particles. It is the thermodynamic cost of the system's transition from
unbound potential to bound actuality --- the heat $Q$ released
irreversibly at the virialization boundary, encoded on the nearest
available surface at Landauer cost $k_B T \ln 2$ per bit.
$\langle K \rangle$ is retained in the system as the kinetic energy
of the bound particles. $Q$ is encoded permanently as entropy at the
boundary. Both are exactly half of $|\langle V \rangle|$ --- one half
paying the cost of existence, one half sustaining existence.

\subsection{Why the $1/2$ is No Longer Unexplained}

The ratio $\langle K \rangle / |\langle V \rangle| = 1/2$ arises
because the $1/r$ potential is homogeneous of degree $-1$. Euler's
theorem gives the mechanical result $2\langle K \rangle = |\langle V
\rangle|$. IAM's Law gives the thermodynamic result: the system
partitions $|\langle V \rangle|$ exactly in half because that is the
unique partition at which the Landauer cost of the transition equals the
kinetic energy of the resulting bound state. Any other partition would
violate either the virial equilibrium condition or the Landauer cost of
the transition. The $1/2$ is not a coincidence. It is a necessity.

\subsection{Scale Invariance: 37 Orders of Magnitude}

Because the virial temperature cancels in Step~3, the identity
$\langle K \rangle = E_\mathrm{Landauer}$ holds independently of
temperature, mass, or size. Table~\ref{tab:virial} shows the
verification across 37 orders of magnitude in physical scale.

\begin{table}[h]
\centering
\small
\caption{The virial $1/2$ partition verified across 37 orders of
magnitude. The temperature cancellation in Eq.~\eqref{eq:T_cancels}
guarantees scale-invariance. The cosmological entry is the strongest
quantitative test: 17 independent MCMC chains return
$\beta_m = 0.1583 \pm 0.0033$ against the virial prediction $0.1575$
at $0.2\sigma$, with zero free parameters.}
\label{tab:virial}
\begin{tabular}{p{4.0cm}p{1.6cm}p{2.0cm}p{4.2cm}}
\toprule
\textbf{System} & \textbf{Scale} & \textbf{Force} &
\textbf{Verification} \\
\midrule
Electron rest mass & $\lambda_C \sim 10^{-13}$~m & EM/Gravity &
  $m_e$ to 6.6~ppm, zero free parameters~\citep{IAM_archive} \\
Atoms (H--Xe, 20 elements) & $10^{-10}$~m & Coulomb &
  Exact, machine precision~\citep{Clementi1974} \\
Molecules (H$_2$--CO$_2$) & $10^{-9}$~m & Coulomb &
  Exact at equilibrium geometry \\
Solar structure & $10^{9}$~m & Gravity & $\sim$10\%~\citep{Christensen2002} \\
White dwarfs (Chandrasekhar) & $10^{7}$~m & Gravity &
  Exact, $1.44\,M_\odot$ observed~\citep{Chandrasekhar1931} \\
Galaxy clusters ($\sim$20) & $10^{23}$~m & Gravity &
  $\sim$20\%~\citep{Evrard1996} \\
N-body virial efficiency & $10^{22}$--$10^{25}$~m & Gravity &
  $0.815 \pm 0.025$, 6 independent sims~\citep{BryanNorman1998} \\
\rowcolor{mcmcrow}
Cosmological $\beta_m$ & $10^{26}$~m & Gravity &
  $0.1583 \pm 0.0033$, 0.2$\sigma$, 17 chains~\citep{IAM_archive} \\
\bottomrule
\multicolumn{4}{l}{\textit{Scale range: $\lambda_C \sim 10^{-13}$~m to $10^{26}$~m
--- 37 orders of magnitude.}} \\
\end{tabular}
\end{table}

% ─────────────────────────────────────────────────────────────────────────────
\section{Global Expression: The Jacobson--Cai-Kim Chain Completed}
\label{sec:global}
% ─────────────────────────────────────────────────────────────────────────────

\subsection{What Einstein and Jacobson Were Missing}

Einstein's field equations describe the geometry of actual spacetime.
They are correct and complete as a description of the geometry that
results from matter distribution. But GR has no account of the
thermodynamic cost of actualization itself. It describes what happens
\textit{after} decoherence events produce classical matter. It is
silent on what it costs the universe to produce that classical matter.

Jacobson~\citep{Jacobson1995} came the closest. He showed that the
Einstein equations follow from the Clausius relation $\delta Q = T\,dS$
applied to local causal horizons. The derivation proceeds as follows.
At any spacetime point $p$, a small spacelike 2-surface $\mathcal{P}$
defines a local Rindler horizon $\mathcal{H}$ with null tangent $k^a$
and approximate boost Killing vector $\chi^a = -\kappa\lambda k^a$,
where $\kappa$ is the surface gravity. The Unruh temperature is
$T = \hbar\kappa/(2\pi k_B)$~\citep{Unruh1976}. The heat flux across
the horizon is:
\begin{equation}
\delta Q = -\kappa \int_\mathcal{H}
\lambda\,T_{ab}\,k^a k^b\,d\lambda\,dA.
\label{eq:heat_flux}
\end{equation}
With entropy proportional to horizon area $S = \eta A$, the entropy
variation via the Raychaudhuri equation gives:
\begin{equation}
\delta S = -\eta \int_\mathcal{H}
\lambda\,R_{ab}\,k^a k^b\,d\lambda\,dA.
\label{eq:dS_jacobson}
\end{equation}
Setting $\delta Q = T\,\delta S$ and requiring this to hold for
\textit{all} null vectors $k^a$ yields the Einstein field equation
with $G = 1/(4\hbar\eta)$~\citep{Jacobson1995}. The machine is:
\textit{specify an entropy functional, derive a gravity theory.}

\subsubsection*{Derivation of $S \propto A$ and $\eta = 1/(4\ell_P^2)$}

In Jacobson's original derivation, both $S \propto A$ and the value of
$\eta$ are treated as inputs. Both are in fact geometric consequences of
flat spacetime.

\textbf{(A) $S \propto A$ from causal accessibility.}
Every irreversible quantum-to-classical transition produces classical
information that must be permanently encoded somewhere accessible to
external observers. Information inside the horizon is causally
inaccessible by definition. The causal boundary is a 2D surface. Therefore
information must be encoded on the boundary, not in the bulk volume.
Information capacity of a 2D surface scales with area: $S \propto A$.
This is not a postulate --- it follows from the irreversibility of
decoherence and the causal structure of the horizon.

\textbf{(B) $\eta = 1/(4\ell_P^2)$ from $2\pi/8\pi$.}
The structural equation from Jacobson's derivation is:
\begin{equation}
\frac{\hbar\eta}{2\pi} = \frac{c^4}{8\pi G}.
\label{eq:jacobson_structure_law}
\end{equation}
The $2\pi$ on the left is the Euclidean Rindler cone angle --- the unique
topological period eliminating the conical singularity at the horizon in flat
spacetime. Under Euclidean continuation $t\to -i\tau$, the Rindler metric
$ds^2 = -\kappa^2\rho^2 dt^2 + d\rho^2 + d\mathbf{x}_\perp^2$ becomes a
cone in the $(\rho, \kappa\tau)$ plane. Regularity at $\rho = 0$ requires
$\kappa\tau \sim \kappa\tau + 2\pi$ --- a topological necessity of flat
spacetime, independent of $\eta$, $S \propto A$, or $G$.

The $8\pi$ on the right is the Einstein normalization: $8\pi = 2 \times 4\pi$,
where $4\pi$ is the solid angle of the unit sphere in $\mathbb{R}^3$ (Gauss's
law), and the factor 2 is the trace structure of $G_{ab} = R_{ab} -
\frac{1}{2}Rg_{ab}$ in the weak-field limit~\citep{Wald1984}. Both factors
are geometric properties of three-dimensional space, independent of $\eta$
or $S \propto A$.

Solving Eq.~\eqref{eq:jacobson_structure_law} for $\eta$:
\begin{equation}
\boxed{\eta = \frac{c^4}{4\hbar G} = \frac{1}{4\ell_P^2},
\qquad \frac{1}{4} = \frac{2\pi}{8\pi}.}
\label{eq:eta_law}
\end{equation}
No observational input is required. Both factors were already present in
Jacobson's derivation. Their ratio fixes $\eta$ as a geometric consequence.

\textbf{(C) Consistency: one bit per $4\ell_P^2$.}
The Landauer cost of one decoherence event at horizon temperature $T_H$ is
$E_\mathrm{bit} = \hbar\kappa/(2\pi)$. The minimum horizon area disturbed:
\begin{equation}
\delta A_\mathrm{min} = \frac{8\pi G}{c^4}\cdot\frac{\hbar\kappa}{2\pi}
= \frac{4\hbar G\kappa}{c^4}
\;\xrightarrow{\kappa = c^2/\ell_P}\;
4\ell_P^2.
\end{equation}
Therefore $\eta \cdot \delta A_\mathrm{min} = 1$: exactly one bit per
$4\ell_P^2$. The entropy coefficient and the minimum encoding area are
the same geometric ratio $8\pi/2\pi = 4$ stated twice. Jacobson's machine
now fully specifies its own entropy functional.

 It contained only $S_\mathrm{geo}
= \eta A$. The Landauer cost of irreversible bulk information production
--- the cost paid each time a quantum system virializes and crosses from
potential to actual --- was not included. IAM's Law identifies this as
the missing term.

Cai and Kim~\citep{CaiKim2005} extended Jacobson's argument to the FRW
apparent horizon, recovering the Friedmann equations from the first
law applied to the apparent horizon at $\tilde{r}_A = 1/H$ with
area $A_H = 4\pi/H^2$, geometric entropy $S_\mathrm{geo} =
\pi/(GH^2)$, and Gibbons-Hawking temperature $T_H = \hbar H/(2\pi
k_B)$~\citep{Gibbons1977}. The same incompleteness applies: the
entropy functional lacked the informational term.

IAM's Law identifies the missing term. The total entropy of the
cosmic horizon has two components:
\begin{equation}
S_\mathrm{total} = S_\mathrm{geo} + S_\mathrm{info},
\label{eq:Stotal}
\end{equation}
where $S_\mathrm{geo} = A_H/4G$ is the Bekenstein-Hawking geometric
entropy~\citep{Bekenstein1973,Hawking1975} that reproduces GR through
the Jacobson mechanism, and $S_\mathrm{info}$ is the accumulated
Landauer cost of every irreversible quantum-to-classical transition
in the bulk since stable massive particles first existed.

\subsection{Two Sources of Information}

The total information production rate in the universe has two
components:

\textbf{Vacuum baseline ($i_\mathrm{vac}$):} Quantum vacuum
fluctuations produce a constant, time-independent rate of information.
This baseline is always present, independent of structure. It maps
directly to the cosmological constant $\Lambda$ through the Jacobson
mechanism.

\textbf{Structural complexity ($i_\mathrm{struct}(a)$):} Gravitational
collapse, decoherence, and bound-state formation produce a
time-dependent growing contribution. In the early universe, before
significant structure exists, $i_\mathrm{struct} \approx 0$. As
galaxies form and the cosmic web develops, $i_\mathrm{struct}$ grows.
As expansion dilutes matter and weakens gravitational collapse, the
process self-regulates toward a finite asymptote. This term maps to
$\beta_m\,\mathcal{E}(a)$ in the matter-sector perturbation equations.

\subsection{The Cai--Kim Derivation and the IAM Modification}

Setting $-dE = T_H\,dS_\mathrm{geo}$ with the Misner-Sharp energy
$E = \tilde{r}_A/2G$ and using
$dS_\mathrm{geo} = -(2\pi/GH^3)\,dH$ recovers the Friedmann
equations~\citep{CaiKim2005}. Adding $S_\mathrm{info}$:
\begin{equation}
-dE = T_H\,d\bigl(S_\mathrm{geo} + S_\mathrm{info}\bigr).
\label{eq:modified_first_law}
\end{equation}
The geometric piece yields the standard Friedmann equations. The
informational piece introduces an additional effective energy density.
\textit{A critical clarification before writing this explicitly: the
term $\beta_m\,\mathcal{E}(a)$ is sourced by inhomogeneous
gravitational structure formation. In an exact homogeneous FRW
spacetime, structure formation vanishes and with it this term.
Equation~\eqref{eq:modified_friedmann} below is the formal
thermodynamic completion of the Jacobson--Cai-Kim entropy balance ---
it is not a universal homogeneous background modification applicable
to all sectors. The observable consequences of this term enter
exclusively through the timelike matter-sector perturbation equations,
as derived in Section~\ref{sec:perturbations}. The background FRW
expansion and all photon-sector observables are unchanged.}
\begin{equation}
H^2 = \frac{8\pi G}{3}\,\rho_m + \frac{\Lambda}{3}
  + \beta_m\,\mathcal{E}(a)\,H_0^2.
\label{eq:modified_friedmann}
\end{equation}
This is not a modification of GR. It is GR's thermodynamic derivation
completed with the full entropy budget.

\subsection{Information Production Rate}

The rate of irreversible information production from gravitational
structure formation scales as:
\begin{equation}
\dot{I} \propto \rho_m(a)\cdot D(a)^{5/2}\cdot f(a)\cdot H(a),
\label{eq:Idot}
\end{equation}
where $\rho_m \propto a^{-3}$ is the matter density, $D(a)$ is the
linear growth factor, $f = d\ln D/d\ln a$ is the growth rate, and
$n = 5/2$ is the critical exponent derived in
Section~\ref{sec:bridge}. The effective rate of informational entropy
increase on the horizon is:
\begin{equation}
\frac{dS_\mathrm{info}}{dt}
= \frac{\dot{I}}{T_H \cdot A_H},
\label{eq:dSinfo_rate}
\end{equation}
where $1/T_H = 2\pi k_B/(\hbar H)$ reflects the Landauer encoding
cost per bit at the horizon temperature, and $1/A_H$ converts total
information to surface density consistent with the holographic
bound~\citep{tHooft1993,Susskind1995}.

\subsection{The Activation Function Derived}

The decoherence constraint requires the informational energy density
to satisfy:
\begin{equation}
\dot{\rho}_\mathrm{info} = \rho_\mathrm{info}\cdot\frac{H}{a}.
\label{eq:rhodot}
\end{equation}
This separable ODE integrates directly:
\begin{equation}
\rho_\mathrm{info}(a) \propto
\exp\!\left(\int_0^a \frac{da'}{a'^2}\right)
= \exp\!\left(1 - \frac{1}{a}\right) \equiv \mathcal{E}(a).
\label{eq:Ea_derived}
\end{equation}
Normalization $\mathcal{E}(1) = 1$ is set by definition. Therefore:
\begin{equation}
\boxed{\mathcal{E}(a) = \exp\!\left(1 - \frac{1}{a}\right).}
\label{eq:Ea}
\end{equation}
The properties of $\mathcal{E}(a)$ follow directly from IAM's Law:
\begin{itemize}
\item $\mathcal{E}(a) \to 0$ as $a \to 0$: no actualization cost
  before stable massive particles exist.
\item $\mathcal{E}(1) = 1$: normalization today.
\item $d\mathcal{E}/da > 0$ always: each decoherence event is
  irreversible. The ledger only grows.
\item $\mathcal{E}(\infty) = e \approx 2.718$: the accumulated
  cost asymptotes without arriving. The universe matures rather
  than tears apart.
\item $\mathcal{E}(z) = \exp(-z)$: pure exponential decay with
  redshift --- the simplest possible form for a quantity growing
  from zero at early times to a finite value today.
\end{itemize}

The $1/a$ in the exponent has a precise physical meaning:
$\mathcal{E}(a)$ is the exponential of the cumulative
\textit{information surface density} on the cosmic horizon. What
drives expansion is not how much total information the universe has
produced, but how densely that information is packed on the horizon
surface.

\subsection{The Variational Formulation}

The thermodynamic result admits a variational formulation, connecting
IAM's Law to standard field-theoretic language. Define the scalar
field~\citep{IAM_archive}:
\begin{equation}
\varphi(t) \equiv \ln\mathcal{E}(a) = 1 - \frac{1}{a},
\label{eq:phi_def}
\end{equation}
which satisfies the decoherence constraint $\dot{\varphi} = H/a$.
This is not a Klein-Gordon equation. $\varphi$ does not propagate
independently; it tracks the cumulative decoherence history. The
informational action is:
\begin{equation}
S_\mathrm{info} = \int\!\left[
  -\frac{3H_0^2}{8\pi G}\,\beta_m\,e^{\varphi}
  + \lambda\!\left(\dot{\varphi} - \frac{H}{a}\right)
\right]\!\sqrt{-g}\;d^4x,
\label{eq:Sinfo}
\end{equation}
where $\lambda$ is a Lagrange multiplier enforcing the decoherence
constraint. Variation with respect to $g_{ab}$ recovers
Eq.~\eqref{eq:modified_friedmann} exactly. The equation of state
of the informational energy density follows from the continuity
equation $\dot{\rho} + 3H(\rho + P) = 0$:
\begin{equation}
w_\mathrm{info}(a) = -1 - \frac{1}{3a}.
\label{eq:w_info}
\end{equation}
This is a mildly phantom equation of state ($w < -1$) at all finite
$a$, approaching $-1$ as $a \to \infty$. The phantom behavior
reflects the continuous conversion of gravitational potential energy
into informational entropy --- not a violation of the null energy
condition, because $\varphi$ is a thermodynamic variable constrained
by geometry, not a propagating field. $w_\mathrm{info}(a)$ is a
formal thermodynamic parametrization of the entropy completion; it
is not a homogeneous dark-energy fluid implemented in the Boltzmann
equations. The operational implementation of IAM in observable
cosmology is through the effective matter-sector friction history
described in Section~\ref{sec:perturbations}, not through a literal
new background source term.

Energy-momentum is conserved identically:
$\nabla_\mu T^{\mu\nu}_\mathrm{total} = 0$. No energy is created or
destroyed. The formal thermodynamic representation
Eq.~\eqref{eq:modified_friedmann} is consistent with the contracted
Bianchi identity and the Einstein tensor; this consistency confirms
the entropy-balance derivation is self-coherent, not that a new
universal homogeneous fluid has been added to the Friedmann equations.

\subsection{The Virial Theorem as Balance Sheet}

IAM's Law tells you that actualization events carry a Landauer cost.
The virial theorem tells you exactly how much energy crosses the
actualization boundary at each virialization event. Together they
give the cosmological coupling constant with zero free parameters.

For a self-gravitating system in virial equilibrium, half the
gravitational binding energy goes into the geometric channel
(spacetime curvature, standard GR) and half goes into the
informational channel (the actualization cost, paid to the encoding
surface). The informational energy density at $a = 1$ is therefore:
\begin{equation}
\rho_\mathrm{info}(a{=}1) = \frac{1}{2}\,\rho_m
= \frac{\Omega_m}{2}\,\rho_\mathrm{crit}.
\end{equation}
Since $\rho_\mathrm{info}(1) = \beta_m\,\mathcal{E}(1)\,\rho_\mathrm{crit}
= \beta_m\,\rho_\mathrm{crit}$:
\begin{equation}
\boxed{\beta_m = \frac{\Omega_m}{2}.}
\label{eq:beta_derived}
\end{equation}
This is derived, not fitted. The $1/2$ partition of the virial theorem
IS the coupling constant of IAM's Law at cosmological scales. With
$\Omega_m = 0.315$~\citep{PlanckVI2020}, the prediction is
$\beta_m = 0.1575$.

\subsection{The Dual-Sector Split}

IAM's Law applies exclusively to matter on timelike worldlines.
\begin{definition}[Decoherence eligibility]
A degree of freedom undergoes gravitational decoherence if and only if
it propagates on a timelike worldline ($d\tau > 0$).
\end{definition}

Decoherence is a dynamical process requiring a finite proper time
interval $\Delta\tau > 0$. Photons propagate on null geodesics:
$ds^2 = 0$, $\Delta\tau = 0$. A photon does not accumulate proper
time. It does not virialize. It pays no actualization cost. This is
a geometric consequence of spacetime structure, not an assumption.

\begin{proposition}[Dual-sector coupling]
\label{prop:dual}
For matter (timelike worldlines): $S_\mathrm{info} > 0$, yielding:
\begin{equation}
\mu(a) = \frac{H^2_{\Lambda\mathrm{CDM}}(a)}
{H^2_{\Lambda\mathrm{CDM}}(a) + \beta_m\,\mathcal{E}(a)} < 1.
\label{eq:mu_exact}
\end{equation}
For photons (null worldlines): $S_\mathrm{info} = 0$, yielding:
\begin{equation}
\Sigma(a) = 1 \quad (\text{exact, geometric necessity}).
\label{eq:sigma_exact}
\end{equation}
\end{proposition}

At $a = 1$: $\mu_0 \equiv \mu(z{=}0) - 1 = -\beta_m/(1+\beta_m)
= -0.136$, with zero free parameters beyond $\Lambda$CDM.

The $\beta_m\,\mathcal{E}(a)$ term does not modify the background
Friedmann equation. In an exact homogeneous FRW spacetime the Weyl
tensor vanishes ($C_{\mu\nu\rho\sigma} = 0$), and any entropy
functional sourced by inhomogeneous structure formation must vanish
in this limit. The IAM modification operates exclusively at the
perturbation level, entering through the effective Hubble friction
for matter perturbations while leaving the background and photon
sector unchanged.

% ─────────────────────────────────────────────────────────────────────────────
\section{Perturbation Theory: Deriving $\mu(a) < 1$ and $\Sigma(a) = 1$}
\label{sec:perturbations}
% ─────────────────────────────────────────────────────────────────────────────

The entropy completion of Section~\ref{sec:global} is a formal
thermodynamic statement. This section shows precisely how it enters
observable cosmology: not as a universal homogeneous background
modification, but exclusively through the timelike matter-sector
perturbation equations. The background FRW expansion rate, photon
geodesics, CMB power spectra, BAO positions, and lensing are all
unchanged. The sole observable effect is suppressed matter growth via
enhanced Hubble friction in the matter-sector growth equation.

\subsection{The Informational Field Does Not Fluctuate}

The scalar field $\varphi = 1 - 1/a$ is defined through the background
apparent horizon radius $\tilde{r}_A = 1/H(t)$ and therefore depends
only on the homogeneous expansion rate $H(t)$. In linear scalar
perturbation theory, local overdensities $\delta\rho$ source metric
potentials $\Psi$ and $\Phi$, but they do not shift the
background-defined apparent horizon at first order. Therefore:
\begin{equation}
\delta\varphi = 0
\label{eq:dphi_zero}
\end{equation}
exactly, at all orders in linear perturbation theory. This is
analogous to the cosmological constant, which satisfies
$\delta\Lambda = 0$: both are properties of the spacetime geometry,
not of the local matter content. Formally, perturbing the constrained
action (Eq.~\eqref{eq:Sinfo}): variation with respect to $\lambda$
gives the perturbed constraint $\delta\dot{\varphi} = 0$, and combined
with the initial condition $\delta\varphi(a_i) = 0$, the field remains
unperturbed at all times.

\subsection{Standard GR Perturbations with Modified Matter-Sector Friction}

With $\delta\varphi = 0$, the cosmological perturbation equations are
identical to standard GR. The sole effect of the entropy completion
is to introduce an additional contribution to the effective
matter-sector expansion rate, denoted $H_\mathrm{eff,m}$, which
enters only the matter growth equation as enhanced Hubble friction.
The background FRW expansion rate is unchanged. In conformal
Newtonian gauge
($ds^2 = a^2[-(1+2\Psi)d\tau^2 + (1-2\Phi)d\mathbf{x}^2]$):

\paragraph{Poisson equation:}
\begin{equation}
k^2\Psi = -4\pi G\,a^2\,\rho_m\,\delta_m.
\label{eq:poisson_iam}
\end{equation}

\paragraph{No anisotropic stress:}
\begin{equation}
\Psi = \Phi
\label{eq:no_stress}
\end{equation}
since $\delta\varphi = 0$ introduces no anisotropic stress from the
informational sector.

\paragraph{Matter growth equation:}
\begin{equation}
\ddot{\delta}_m + 2H_\mathrm{eff,m}\,\dot{\delta}_m
- 4\pi G\,\rho_m\,\delta_m = 0,
\label{eq:growth_IAM}
\end{equation}
where $H_\mathrm{eff,m}^2 \equiv H_{\Lambda\mathrm{CDM}}^2
+ \beta_m\,\mathcal{E}(a)\,H_0^2$ is the effective matter-sector
friction history derived from the entropy completion.
These are standard GR perturbation equations. The only modification
from $\Lambda$CDM is the enhanced Hubble friction term
$2H_\mathrm{eff,m}\,\dot{\delta}_m$, which suppresses matter growth.
This is the mechanism of $\mu(a) < 1$.

\subsection{$\mu(a)$ and $\Sigma(a)$ as Ratios of Hubble Rates}

Mapping to the $\mu$--$\Sigma$ framework, the effective gravitational
coupling for matter perturbations is the ratio of Hubble rates squared:
\begin{equation}
\boxed{\mu(a) = \frac{E^2_{\Lambda\mathrm{CDM}}(a)}{E^2_\mathrm{eff,m}(a)}
= \frac{H^2_{\Lambda\mathrm{CDM}}(a)}
{H^2_{\Lambda\mathrm{CDM}}(a) + \beta_m\,\mathcal{E}(a)\,H_0^2} < 1.}
\label{eq:mu_derived}
\end{equation}
The suppression $\mu < 1$ arises because the enhanced Hubble friction
in IAM reduces the effective matter density parameter:
$\Omega_m^\mathrm{IAM}(a) < \Omega_m^{\Lambda\mathrm{CDM}}(a)$,
so matter clusters less effectively than in $\Lambda$CDM. Numerical
values at key redshifts: $\mu(z{=}0) = 0.864$, $\mu(z{=}0.5) = 0.948$,
$\mu(z{=}1) = 0.982$, with $\mu \to 1$ at high redshift as
$\mathcal{E}(a) \to 0$.

Since both potentials satisfy the unmodified Poisson equation
(Eq.~\eqref{eq:poisson_iam}) with no anisotropic stress
(Eq.~\eqref{eq:no_stress}), photon geodesics are determined by
$(\Psi + \Phi)/2 = \Psi$ without modification. Therefore:
\begin{equation}
\Sigma(a) = 1 \quad (\text{exact}).
\label{eq:sigma_derived}
\end{equation}
The photon exemption at the perturbation level is a \textit{consequence}
of $\delta\varphi = 0$ --- that $\varphi$ is a horizon quantity, not a
local field --- not an additional assumption. CMB power spectra, CMB
lensing, BAO angular positions, and supernova luminosity distances are
unchanged. IAM and $\Lambda$CDM produce CMB TT spectra differing by
less than $0.13\%$ at $\ell > 30$~\citep{IAM_archive}.

\subsection{Uniqueness in the $\mu$--$\Sigma$ Parameter Space}

The combination $\mu(a) < 1$ with $\Sigma(a) = 1$ is unique among
known theoretical models. Most modified gravity theories predict
$\mu > 1$ (enhanced growth) or correlated deviations in both $\mu$
and $\Sigma$:
\begin{itemize}
\item $f(R)$ gravity: $\mu > 1$, $\Sigma > 1$
\item DGP braneworld: $\mu > 1$, $\Sigma \approx 1$
\item General Horndeski: correlated deviations in both $\mu$ and $\Sigma$
\item IAM: $\mu < 1$, $\Sigma = 1$ exactly
\end{itemize}
The IAM signature --- weakened matter growth with unmodified lensing
--- is characteristic of a model that modifies matter-sector dynamics
through Hubble friction without introducing anisotropic stress or
touching photon geodesics. No other known framework produces this
combination from first principles.

\subsection{Numerical Validation}

Level~1 validation uses MGCAMB v1.5.2 + Cobaya in the $\mu$--$\Sigma$
parametrization (12 chains, 4 dataset combinations: Planck only,
Planck + RSD, Planck + BAO, Planck + Pantheon+). Level~2 validation
uses direct Fortran modification of CAMB v1.5.8 + Cobaya implementing
the IAM growth equation (Eq.~\eqref{eq:growth_IAM}) via the matter-sector
expansion rate \texttt{adotoa\_matter} in the Boltzmann solver (5 chains).
Two Level~2b exploratory runs applying the modification at the background
Friedmann level produced $H_0 \approx 61.5$~km\,s$^{-1}$\,Mpc$^{-1}$,
confirming that IAM operates exclusively at the perturbation level.
All 17 chains converged to $R-1 < 0.01$ with zero discarded
runs~\citep{IAM_archive}.

% ─────────────────────────────────────────────────────────────────────────────
\section{The Bridge: $n = 5/2$ from Local to Global}
\label{sec:bridge}
% ─────────────────────────────────────────────────────────────────────────────

The claim that local and global are the same law requires a
quantitative bridge. That bridge is the critical exponent $n = 5/2$,
derived independently from two directions.

\subsection{Top-Down: From Horizon Thermodynamic Consistency}

The top-down derivation requires the cumulative informational entropy
to reproduce $\mathcal{E}(a) = \exp(1-1/a)$. Under matter-domination
scalings ($H \propto a^{-3/2}$, $A_H \propto a^3$, $T_H \propto
a^{-3/2}$, $D \propto a$), the integrand of
Eq.~\eqref{eq:dSinfo_rate} per unit $\ln a$ gives:
\begin{equation}
\frac{dS_\mathrm{info}}{d\ln a} \propto
\frac{\rho_m \cdot D^n \cdot f}{T_H \cdot A_H}
\propto a^{n - 9/2}.
\label{eq:topdown_scaling}
\end{equation}
Converting to $da$ and integrating:
\begin{equation}
S_\mathrm{info}(a) \propto \frac{a^{n-9/2}}{n - 9/2}.
\label{eq:Sinfo_integral}
\end{equation}
For $S_\mathrm{info}(a) \propto -1/a + \text{const}$ (the form
required by $\mathcal{E}(a) = \exp(1-1/a)$), we need
$a^{n-9/2} = a^{-1}$:
\begin{equation}
n - \frac{9}{2} = -1
\quad\Longrightarrow\quad \boxed{n = \frac{5}{2}.}
\label{eq:n_topdown}
\end{equation}

\subsection{Bottom-Up: From Structure Formation Physics}

The bottom-up derivation integrates IAM's Law over the Press-Schechter
halo mass function~\citep{PressSchechter1974}. The rate of
information production per virialization event is set by the Landauer
cost at the virialization boundary (local IAM's Law). Integrating over
the halo population with peak-height scaling $\nu \propto
D(a)^{-1/2}$ near the characteristic mass $M_*$, the dominant
contribution to the information production rate scales as $D^{5/2}$,
giving $n = 5/2$ from quantum decoherence physics directly.

The Sheth-Tormen mass function~\citep{ShethTormen1999} at
$\sigma_* = 1.2$ (galaxy-scale halos where the bulk of cosmic
information production occurs) recovers $\mathcal{E}(a) \propto
\exp(0.925 - 1.009/a)$ with Pearson correlation $r = 0.992$ ---
the $1/a$ exponent to within $1\%$ without free parameters.
Four independent N-body studies spanning
$z = 0$--$30$~\citep{Jenkins2001,Reed2003,Tinker2008,Watson2013}
measure $n_\mathrm{eff} = 3.33 \pm 0.33$, consistent with the IAM
predicted range $n \approx 3.0$--$3.5$ under full $\Lambda$CDM
evolution beyond the matter-domination approximation.

\subsection{Convergence: One Law, Two Directions}

Both directions yield $n = 5/2$. The top-down derivation requires
it from horizon thermodynamic consistency. The bottom-up derivation
derives it from the physics of quantum decoherence integrated over
the halo population. Their agreement is a nontrivial consistency
result: the same thermodynamic cost per bit, applied at the local
decoherence boundary and integrated across the halo population,
recovers the same critical exponent that horizon thermodynamics
independently requires. This is strong support for one mechanism
operating at two scales.

% ─────────────────────────────────────────────────────────────────────────────
\section{The Nine-Step Causal Chain}
\label{sec:chain}
% ─────────────────────────────────────────────────────────────────────────────

This section traces the causal sequence by which IAM's Law, applied
to gravitational structure formation, produces the observable
cosmological consequences derived in Sections~\ref{sec:global}
and~\ref{sec:perturbations}. Each step follows from established
physics; the references are to the standard results used at each
stage.

\begin{enumerate}
\item \textbf{Gravitational interaction.} Matter interacts
  gravitationally, initiating collapse of overdense regions
  (GR~\citep{Einstein1915}).

\item \textbf{Gravitational collapse.} Matter clusters into bound
  systems: halos, galaxies, stars
  (structure formation~\citep{Peebles1980}).

\item \textbf{Quantum decoherence.} Gravitational interaction
  resolves quantum superpositions into definite classical states
  (Eq.~\eqref{eq:decoherence},~\citep{Zurek2003}).

\item \textbf{Irreversible information production.} Each
  decoherence event produces classical information ($I =
  \log_2 N$ bits, Eq.~\eqref{eq:info_produced}) that cannot be
  recovered.

\item \textbf{Landauer thermodynamic cost.} Information production
  carries an irreducible energy cost of $k_B T_H \ln 2$ per bit
  at the horizon temperature (IAM's Law, Eq.~\eqref{eq:iams_law}).

\item \textbf{Holographic encoding.} The produced information is
  bounded by and encoded on the cosmic horizon
  area~\citep{Bekenstein1973,tHooft1993,Susskind1995},
  contributing $S_\mathrm{info}$ to the entropy functional.

\item \textbf{Informational pressure.} The modified first law
  $-dE = T_H\,d(S_\mathrm{geo} + S_\mathrm{info})$ introduces
  additional effective energy density $\beta_m\,\mathcal{E}(a)H_0^2$
  in the matter-sector perturbation equations
  (Eq.~\eqref{eq:modified_friedmann}).

\item \textbf{Modified growth.} The matter-sector perturbation
  equations acquire enhanced Hubble friction, suppressing growth:
  $\mu(a) = H^2_{\Lambda\mathrm{CDM}}/H^2_\mathrm{IAM} < 1$.

\item \textbf{Feedback suppression.} Suppressed growth dilutes
  $\Omega_m(a)$, weakening gravitational collapse and reducing
  the amplitude of subsequent decoherence events.
\end{enumerate}

Step~9 feeds back into Step~1 with reduced amplitude. The chain
self-regulates, and $\mathcal{E}(a)$ asymptotes to $e$ at late
times rather than diverging. This is the mechanism by which the
activation function acquires its finite upper bound.

Note that gravity appears at two points in the chain: as the
initiating cause of decoherence (Step~1) and as the observable
modified by the resulting informational pressure (Step~8). The
$\beta_m = \Omega_m/2$ posterior in Section~\ref{sec:obs},
consistent with the virial theorem prediction at $0.2\sigma$,
is the quantitative signature of this feedback structure.

% ─────────────────────────────────────────────────────────────────────────────
\section{Black Holes: The Local Saturation Limit of IAM's Law}
\label{sec:bh}
% ─────────────────────────────────────────────────────────────────────────────

The results of Sections~\ref{sec:local}--\ref{sec:perturbations}
address the distributed regime of IAM's Law: structure formation
producing decoherence events across many halos, encoding
cumulatively on the growing cosmic horizon. This section addresses
the concentrated local limit: what IAM's Law predicts when the
local information production rate saturates the encoding capacity
of the nearest bounding surface. The result is the standard black
hole formation condition, recovered here from the informational
saturation criterion rather than from gravitational collapse
directly.

\subsection{The Holographic Saturation Condition}

The feedback in the nine-step chain becomes runaway when the local
information production rate saturates the encoding capacity of the
bounding surface. A black hole forms when the informational entropy
from gravitational decoherence within radius $R$ saturates the
Bekenstein-Hawking bound on the bounding surface:
\begin{equation}
\frac{S_\mathrm{info}(R)}{A(R)} \geq \frac{1}{4\ell_P^2}.
\label{eq:saturation}
\end{equation}
Setting $S_\mathrm{info} = S_\mathrm{BH} = 4\pi GM^2/(\hbar c)$
at the Schwarzschild radius $R_s = 2GM/c^2$:
\begin{equation}
\frac{4\pi GM^2/(\hbar c)}{4\pi\cdot 4G^2M^2/c^4}
= \frac{c^3}{4\hbar G}
= \frac{1}{4\ell_P^2}.
\end{equation}
The holographic bound is exactly saturated at the Schwarzschild radius.
Eq.~\eqref{eq:saturation} is identically equivalent to the Thorne
hoop conjecture~\citep{Thorne1972}: the hoop conjecture is the
holographic bound in geometric language; the saturation condition is
the hoop conjecture in informational language.

\subsection{The Two-Horizon Framework}

Both the black hole horizon (temperature $T_\mathrm{BH} = \hbar c^3
/(8\pi GMk_B)$) and the cosmic horizon (temperature $T_H = \hbar H/
(2\pi k_B)$~\citep{Gibbons1977}) are thermal encoding surfaces within
the IAM framework. Setting $T_\mathrm{BH} = T_H$ yields the thermal
equilibrium mass:
\begin{equation}
M_\mathrm{eq} = \frac{c^3}{4GH}.
\label{eq:Meq}
\end{equation}
At the present epoch, $M_\mathrm{eq} \approx 2.3 \times 10^{22}
M_\odot$ --- far exceeding the largest known black hole. All
astrophysical black holes satisfy $T_\mathrm{BH} \gg T_H$: information
flows unambiguously from hot local horizons to the cold cosmic
horizon, at the thermally-limited rate $\Gamma = c^3/(1920\,GM\ln 2)$
bits\,s$^{-1}$. The cosmic horizon is the final repository; all
locally-encoded information migrates to the global ledger.

\subsection{Black Holes as Preferred Early-Universe Encoding Surfaces}

The activation function $\mathcal{E}(a) = \exp(1-1/a)$ tracks the
encoding transition quantitatively. At early times ($a \ll 1$), the
cosmic horizon is small and the Gibbons-Hawking temperature is high,
making cosmic horizon encoding inefficient. Information from early
gravitational decoherence is encoded primarily on local black hole
horizons. As the universe expands, large-scale structure formation
produces distributed decoherence, the cosmic horizon grows and cools,
eventually dominating as the encoding surface. The activation function
is the quantitative record of this transition.

This suggests a natural thermodynamic context for the observed
population of massive black holes at high redshift: in the early
universe, before the cosmic horizon grows large enough to encode
efficiently, black holes are the nearest available encoding surfaces.
Within the IAM framework, this is a natural consequence of the
saturation condition rather than an anomaly requiring separate
explanation.

% ─────────────────────────────────────────────────────────────────────────────
\section{The Thermodynamic Arrow of Time and the Coincidence Problem}
\label{sec:arrow}
% ─────────────────────────────────────────────────────────────────────────────

IAM's Law is irreversible by construction: each Landauer cost payment
is permanent, and $S_\mathrm{info}$ can only grow. This section
records two direct consequences of that irreversibility that are
not independent predictions but structural properties of the
activation function $\mathcal{E}(a)$ already derived.

The activation function $\mathcal{E}(a) = \exp(1-1/a)$ is
monotonically increasing in $a$. Since decoherence events are
irreversible and Landauer costs are permanent, $\mathcal{E}(a)$
cannot decrease. Expressed in redshift:
\begin{equation}
\mathcal{E}(z) = \exp(-z).
\label{eq:E_redshift}
\end{equation}
The thermodynamic arrow of time follows directly: the informational
content of the universe increases monotonically with cosmic time.
At recombination ($a \approx 10^{-3}$), $\mathcal{E}(a) \approx 0$;
at the present epoch, $\mathcal{E}(1) = 1$; at late times,
$\mathcal{E}(a) \to e$. The arrow is not an additional assumption.
It is a property of $\mathcal{E}(a)$ inherited from the
irreversibility of decoherence encoded in IAM's Law.

The coincidence problem --- why the informational energy density
$\rho_\mathrm{info} \propto \mathcal{E}(a)$ dominates at
approximately the same epoch as structure formation peaks --- has
a direct causal explanation within IAM. The informational pressure
term $\beta_m\,\mathcal{E}(a)$ is sourced by structure formation.
It grows when structure formation is active and self-regulates as
suppressed growth reduces the decoherence rate. The coincidence
between dark energy domination and the peak of structure formation
is therefore a consequence of IAM's Law rather than a free
parameter: the same process drives both.

% ─────────────────────────────────────────────────────────────────────────────
\section{Observational Consistency}
\label{sec:obs}
% ─────────────────────────────────────────────────────────────────────────────

\subsection{The Primary Test: $\beta_m = \Omega_m/2$}

The primary quantitative test of IAM's Law is the cosmological
coupling constant $\beta_m = \Omega_m/2 = 0.1575$, derived from
the virial theorem with zero free parameters. Seventeen MCMC chains
were run against the full Planck~2018 likelihood (TT, TE, EE +
low-$\ell$ + lensing). All 17 converged to $R-1 < 0.01$ with no
discarded runs, no parameter adjustments, and no excluded redshift
ranges~\citep{IAM_archive}. The $\beta_m$ posterior is consistent
with the prediction of IAM's Law at cosmological scales: the virial
theorem predicted $0.1575$ from first principles; the chains returned
$0.1583 \pm 0.0033$, a $0.2\sigma$ agreement.

\begin{table}[h]
\centering
\small
\caption{Summary of quantitative consistency checks. All results derived
from IAM's Law with zero free parameters beyond
$\Lambda$CDM. The cosmological constant and baryon asymmetry entries
represent quantities derived by IAM's Law that are free parameters in
all other frameworks~\citep{IAM_archive}.}
\label{tab:chains}
\begin{tabular}{p{4.5cm}p{3.6cm}p{4.4cm}}
\toprule
\textbf{Quantity} & \textbf{Prediction} & \textbf{Result} \\
\midrule
$\beta_m$ (virial theorem) &
  $0.1575$ &
  $0.1583 \pm 0.0033$ ($0.2\sigma$) \\
$\sigma_8$ &
  $0.800$ &
  $0.7998 \pm 0.0058$ ($0.1\sigma$) \\
$H_0^\mathrm{(photon)}$ &
  67.16~km\,s$^{-1}$\,Mpc$^{-1}$ &
  $67.16 \pm 0.47$ ($0.37\sigma$ from Planck) \\
$H_0^\mathrm{(matter)}$ &
  72.26~km\,s$^{-1}$\,Mpc$^{-1}$ &
  $72.26$ ($0.75\sigma$ from SH0ES) \\
$\Delta\chi^2$ vs $\Lambda$CDM &
  $\leq 0$ &
  $+0.54$ (best-fit) \\
\rowcolor{newresult}
$\Lambda/\rho_\mathrm{vac}$ (CC) &
  $1.141\times10^{-123}$ &
  $1.14\times10^{-123}$ (0.07\%, 0 free params) \\
\rowcolor{newresult}
$\eta$ (baryon asymmetry) &
  $6.115\times10^{-10}$ &
  $(6.115\pm0.037)\times10^{-10}$ (0.36\%, BBN prior off) \\
\midrule
\multicolumn{3}{l}{All 18 chains converged (17 IAM growth equation;
  1 baryon asymmetry). Zero discarded runs. Zero parameter adjustments.} \\
\multicolumn{3}{l}{\textit{Green rows: quantities derived by IAM's Law
  that are free parameters in all other frameworks.}} \\
\bottomrule
\end{tabular}
\end{table}

\subsection{Independent N-body Confirmation}

Six independent N-body simulation
campaigns~\citep{BryanNorman1998,Bett2007,Neto2007,Ludlow2010,Power2012,Klypin2016}
spanning 25 years measure the mass-weighted virial efficiency
$\langle\eta_\mathrm{vir}\rangle = 0.815 \pm 0.025$, consistent with
the IAM prediction of $0.81$. The combined measurement
$\Omega_m \times f_\mathrm{coll} \times
\langle\eta_\mathrm{vir}\rangle = 0.315 \times 0.62 \times 0.815
= 0.159 \pm 0.010$ agrees with $\beta_m = \Omega_m/2 = 0.1575$
to $1.0\%$ with no free parameters. These simulations were performed
without any knowledge of the IAM framework; their agreement is an
independent cross-check.

\subsection{The Sector-Dependent Hubble Split}

The sector-dependent Hubble split $H_0^\mathrm{(matter)}/
H_0^\mathrm{(photon)} = \sqrt{1 + \beta_m} = 1.073$ is a direct
consequence of Proposition~\ref{prop:dual}. The two values are not
a tension. They are measurements of two distinct physical quantities
separated by the accumulated actualization cost of 13.8 billion
years of structure formation. Both emerge from a single posterior
with zero free parameters beyond $\Lambda$CDM.

Gravitational wave standard sirens provide an independent test.
Binary neutron star mergers occur in galaxies --- matter-sector objects
--- so their host galaxy recession velocities probe the matter-sector
Hubble flow. IAM predicts:
\begin{equation}
H_0^\mathrm{sirens} = H_0^\mathrm{CMB}\sqrt{1 + \beta_m}
= 67.4\sqrt{1.1575} = 72.51\;\mathrm{km\,s^{-1}\,Mpc^{-1}}.
\label{eq:gw_sirens}
\end{equation}
The GW170817 standard siren measurement gives
$H_0 = 75.5^{+5.3}_{-5.4}$~km\,s$^{-1}$\,Mpc$^{-1}$~\citep{Abbott2017},
consistent with the IAM prediction at $0.5\sigma$. This is an
independent consistency check from a completely different observable:
no CMB, no distance ladder, no BAO. The same $\beta_m = \Omega_m/2$
derived from the virial theorem is consistent with the gravitational wave
Hubble constant to within $0.5\sigma$ without adjustment.

\subsection{Precision Measurement of $\mu_0 = -0.136$}

IAM's Law derives $\mu_0 = -0.136$ from first principles with zero
free parameters. Table~\ref{tab:precision} shows the precision with
which upcoming surveys will constrain this quantity. Current data
are already consistent with the prediction at $0.6\sigma$. Tighter
constraints from Euclid and DESI Year~5 will either sharpen the
agreement or reveal a discrepancy --- both outcomes carry information
about the validity of the entropy-completion identification at the
heart of IAM's Law.

\begin{table}[h]
\centering
\small
\caption{Precision with which upcoming surveys will measure the
IAM's Law prediction $\mu_0 = -0.136$. The precision improves by
roughly an order of magnitude from current data to Euclid + DESI
Year~5 combined. Current data are already consistent with the
prediction at $0.6\sigma$~\citep{IAM_archive}.}
\label{tab:precision}
\begin{tabular}{p{4.5cm}p{2.2cm}p{4.5cm}}
\toprule
\textbf{Survey} & \textbf{$\sigma(\mu_0)$} &
\textbf{Precision on $\mu_0 = -0.136$} \\
\midrule
DESI DR1 (current)~\citep{IAM_archive} & 0.22 & $0.6\sigma$ measurement \\
DESI Year 5 (projected) & $\sim 0.08$ & $1.7\sigma$ measurement \\
Euclid DR1 (projected)~\citep{Frusciante2025} & $\sim 0.04$ &
  $3.4\sigma$ measurement \\
Euclid + DESI Year 5 & $\sim 0.03$ & $4.5\sigma$ measurement \\
\bottomrule
\end{tabular}
\end{table}

\subsection{The Cosmological Constant and Baryon Asymmetry as Derived
Quantities}
\label{sec:cc_baryon}

The results reported in this subsection are qualitatively distinct from
those above. The preceding subsections document IAM's Law producing
predictions \textit{consistent with} observational data --- a necessary
but not sufficient condition for a law. This subsection documents
IAM's Law \textit{selecting} two numbers that are free parameters in
every other framework in physics. This is a different category of evidence.

\textbf{The cosmological constant.}
The observed ratio $\Lambda_\mathrm{obs}/\rho_\mathrm{vac} = 1.14
\times 10^{-123}$ has resisted derivation for over fifty years. Within
IAM, this ratio is the fraction of Planck-area horizon bits actualized
by baryonic decoherence over cosmic history. The baseline prediction
\citep{IAM_archive}:
\begin{equation}
\frac{\Lambda_\mathrm{obs}}{\rho_\mathrm{vac}} =
\frac{2}{\pi}\left(\frac{l_P}{l_H}\right)^2
\frac{\Omega_b}{\Omega_m}
= 1.38 \times 10^{-123}
\label{eq:cc_baseline_law}
\end{equation}
recovers the observed value within a factor of 1.2 with zero free
parameters. The factor arises from the departure from exact de~Sitter
geometry. The de~Sitter encoding surface (the cosmological horizon)
has a Gibbons--Hawking temperature $T_{dS} = T_H\sqrt{\Omega_\Lambda}$,
which differs from the effective Hubble writing temperature $T_H$ by
the ratio:
\begin{equation}
\frac{T_{dS}}{T_H} = \sqrt{\Omega_\Lambda} = \sqrt{0.6846} = 0.8274.
\label{eq:tdS_ratio}
\end{equation}
This is a geometric correction, not a free parameter. Applying it:
\begin{equation}
\frac{\Lambda_\mathrm{obs}}{\rho_\mathrm{vac}}\bigg|_\mathrm{corrected}
= \frac{2}{\pi}\left(\frac{l_P}{l_H}\right)^2
\sqrt{\Omega_\Lambda}\,\frac{\Omega_b}{\Omega_m}
= 1.141 \times 10^{-123}
\label{eq:cc_corrected_law}
\end{equation}
against the observed $1.14 \times 10^{-123}$ --- agreement to within
$0.07\%$, zero free parameters beyond the standard cosmological model.
The ratio $\Lambda/\rho_\mathrm{vac}$ is recovered from a calculation
involving only the Planck length, the Hubble constant, and the standard
cosmological density parameters.

\textbf{The baryon asymmetry.}
The baryon-to-photon ratio $\eta \approx 6.14 \times 10^{-10}$ is
independently measured by Big Bang nucleosynthesis and CMB acoustic
peaks, but not derived from a deeper principle in any existing
framework. IAM's Law connects $\eta$ to $\Lambda$ through the
information writing constraint: the CC formula
(Eq.~\ref{eq:cc_baseline_law}) contains $\Omega_b$ explicitly.
Inverting it --- asking what $\Omega_b$ the observed $\Lambda$ requires
--- constrains $\eta$ through $\eta \approx 2.74 \times 10^{-8}
\,\Omega_b h^2$.

A dedicated MCMC chain was run with the BBN prior on $\Omega_b h^2$
removed and replaced with a flat prior four times wider than standard
($[0.010, 0.040]$, compared to the standard $[0.020, 0.025]$). The
chain configuration was otherwise identical to the IAM Level~1 Run~A
used in the primary cosmological test~\citep{IAM_archive}: MGCAMB
v1.5.2 + Cobaya, full Planck~2018 likelihood suite, $\mu_0 = -0.13495$
fixed. The chain has no access to nuclear physics or light element
abundance data. The baryon density is determined solely by the CMB
acoustic peak structure.

The chain converged with $R-1 = 0.009273$ after 22{,}400 accepted
steps. The posterior for $\Omega_b h^2$
yields $\Omega_b h^2 = 0.022319 \pm 0.000136$, corresponding to:
\begin{equation}
\eta_\mathrm{IAM} = (6.115 \pm 0.037) \times 10^{-10}
\label{eq:eta_result_law}
\end{equation}
against the observed $\eta_\mathrm{obs} = 6.137 \times
10^{-10}$~\citep{Cyburt2016,PlanckVI2020} --- agreement to within
$0.36\%$, with no nuclear physics input. Thus the baryon density
selected by the CMB likelihood can be compared directly against the
value required by the IAM cosmological-constant constraint, without
importing BBN as an external prior. The prediction was stated in
advance of the result~\citep{IAM_archive}. No post-hoc adjustment
was made.

\textbf{The connection.}
The cosmological constant problem and the baryon asymmetry problem
are not two independent problems within IAM. They are two expressions
of the same information writing constraint: IAM's Law requires a
specific baryonic matter density to produce the observed $\Lambda$,
and the Planck CMB likelihood, with no external baryon constraint,
returns a density consistent with that requirement. The same
$\sqrt{\Omega_\Lambda}$ geometric correction that closes the CC
calculation to $0.07\%$ also accounts for the offset between the
baseline CC-implied $\eta$ and the CMB-only baryon result, suggesting
a common physical origin within the IAM framework.

These results are documented in \citet{IAM_archive} with timestamps
confirming all predictions were established prior to any chain runs.

% ─────────────────────────────────────────────────────────────────────────────
\section{Assumptions}
\label{sec:assumptions}
% ─────────────────────────────────────────────────────────────────────────────

Any paper claiming to state a law must be explicit about what is
assumed and what is derived. The following assumptions underlie
IAM's Law and its consequences derived in this paper.

\textbf{Assumption 1 --- Gravity is thermodynamic
(Jacobson hypothesis).} The Einstein field equations are an equation
of state derivable from the Clausius relation applied to causal
horizons with entropy proportional to area. This hypothesis is
supported by the Bekenstein-Hawking entropy, the Unruh effect, and
the Cai-Kim derivation of the Friedmann equations from the same
thermodynamic argument. It is the foundation of Jacobson's 1995
result and is not universally accepted.

\textbf{Assumption 2 --- Decoherence-produced information
contributes to horizon entropy.} The Landauer cost of every
irreversible quantum-to-classical transition in the bulk is encoded
on the nearest available encoding surface and contributes to the
entropy functional. This is the single new physical identification
introduced by IAM's Law. It is consistent with the holographic
principle~\citep{tHooft1993,Susskind1995} and Landauer's
principle~\citep{Landauer1961}, but has not been independently
verified at the level of a UV-complete quantum gravity theory. This
is the same gap that exists in all thermodynamic approaches to
gravity.

\textbf{Assumption 3 --- The virial coupling $\beta_m =
\Omega_m/2$ is exact.} The virial theorem in its strict form applies
to systems in stable bound equilibrium. Cosmological structure
formation is non-equilibrium. The derivation in Section~\ref{sec:global}
applies the virial partition to the ensemble of virialized structures
at the population level. The 0.3\% agreement between the virial
prediction and the N-body measurement (Section~\ref{sec:obs}) and the
$0.2\sigma$ MCMC confirmation are empirical justifications, not proofs
of exactness.

\textbf{What is derived, not assumed.} The activation function
$\mathcal{E}(a) = \exp(1-1/a)$ is derived from integrating the
decoherence constraint directly (Section~\ref{sec:global}). The
dual-sector split $\mu < 1$, $\Sigma = 1$ is derived from the
geometric fact that photons propagate on null geodesics
($d\tau = 0$). The critical exponent $n = 5/2$ is derived from both
directions independently (Section~\ref{sec:bridge}). The prediction
$\mu_0 = -0.136$ follows from these inputs with zero free
parameters. None of these are assumed.

\textbf{Three distinct failure modes.} The assumptions above imply
three separable levels at which the framework could be falsified,
and these levels do not rise or fall together.

\textit{Level 1 (core law):} Assumption~2 could fail. Decoherence-produced
information may not contribute to horizon entropy in the way
identified here. Falsification at this level would require a
UV-complete theory that explicitly excludes this channel, or a
derivation showing the entropy identification is internally
inconsistent.

\textit{Level 2 (cosmological implementation):} The virial
coupling Assumption~3 could fail at the population level. A
discrepancy between the predicted $\beta_m = \Omega_m/2$ and
future MCMC posteriors --- or a systematic failure of the
perturbation-level growth suppression to fit upcoming survey data
--- would falsify the cosmological model without necessarily
falsifying the core law.

\textit{Level 3 (broader structural implications):} The
extensions in Sections~\ref{sec:chain}--\ref{sec:arrow} and the
discussion of the weak force and black holes as encoding surfaces
are consequences suggested by the framework but not yet verified
at the same level as the core derivation or the cosmological
implementation. Failure of these extensions would not constitute
falsification of either Level~1 or Level~2.

% ─────────────────────────────────────────────────────────────────────────────
\section{Discussion}
\label{sec:discussion}
% ─────────────────────────────────────────────────────────────────────────────

\subsection{The Three-Step Derivation Chain}

\textbf{Step 1 (Jacobson):} $\delta Q = T\,dS$ on local Rindler
horizons, with $S = \eta A$, implies $G_{ab} + \Lambda g_{ab}
= 8\pi G\,T_{ab}$.

\textbf{Step 2 (Cai--Kim):} $-dE = T\,dS_\mathrm{geo}$ on the FRW
apparent horizon, with $S_\mathrm{geo} = A_H/(4G)$ and
$T = H/(2\pi)$, implies $H^2 = (8\pi G/3)\rho + \Lambda/3$.

\textbf{Step 3 (IAM):} $-dE = T\,d(S_\mathrm{geo} + S_\mathrm{info})$,
with $S_\mathrm{info}$ from gravitational decoherence, implies
$H^2 = (8\pi G/3)\rho + \Lambda/3 + \beta_m\,\mathcal{E}(a)\,H_0^2$.

Step~3 is the only new physical input: gravitational decoherence
irreversibly produces classical information, which is holographically
encoded on the cosmic horizon, contributing $S_\mathrm{info}(a)$ that
grows monotonically with cosmic time. Every other element is
established physics.

\subsection{What Einstein, Jacobson, and Clausius Had}

\textit{Clausius (1870):} The mechanical face of the virial theorem
--- $2\langle K \rangle + \langle V \rangle = 0$ --- derived from
Newton's second law. Correct and complete as a mechanical statement.
Silent on the thermodynamic meaning of $K$.

\textit{Einstein (1915):} The geometry of actual spacetime. Correct
and complete as a description of curved geometry. Silent on the cost
of producing classical matter.

\textit{Jacobson (1995):} Gravity as thermodynamics. GR derived as
an equation of state from $\delta Q = T\,dS$. The machine:
\textit{specify an entropy, derive a gravity theory}. Missing one
entropy term.

\textit{This work:} The missing entropy term is the accumulated
Landauer cost of actualization. Adding it closes Jacobson's
derivation. The thermodynamic face of the virial theorem is the same
cost at the local scale. IAM's Law is the statement that unifies both.

\subsection{The Four Fundamental Forces}

Three of the four fundamental forces satisfy the virial $1/2$
partition. Gravity: confirmed from stellar structure to the
cosmological coupling constant at $0.2\sigma$ across 17 MCMC chains.
Electromagnetism: confirmed at machine precision for 20 atoms and 10
molecules from published Hartree--Fock energies~\citep{Clementi1974},
with mean virial ratio $\eta = 1.0000000000$, standard deviation zero
--- an exact verification because the virial theorem must hold exactly
for converged $1/r$ wavefunctions. The strong nuclear force: the
Chandrasekhar white dwarf limit $M_\mathrm{Ch} = 1.44\,M_\odot$ is
an exact observational consequence of virial equilibrium between
electron degeneracy pressure and gravitational collapse,
confirmed~\citep{Chandrasekhar1931}.

The weak nuclear force does not satisfy the partition. This is
consistent with the physical content of IAM's Law, and the
reason is structural rather than coincidental.

The weak force is excluded by the law's own logic for three
independent reasons. First, it violates parity~\citep{Wu1957}:
the laws governing weak interactions are not symmetric under
spatial reflection. Second, it violates CP symmetry~\citep{Cronin1964}:
it treats matter and antimatter differently, producing the baryon
asymmetry of approximately one excess matter particle per $10^9$
matter-antimatter pairs that constitutes all observable matter.
Third, and most directly: a force whose function is to drive
irreversible transitions \textit{between} stable bound states
cannot simultaneously satisfy an equilibrium condition within them.

The weak force is not one of four equals in the IAM framework.
It is the \textit{precondition} for IAM's Law to operate. Without
CP violation and parity violation, there is no preferred direction
for decoherence, no arrow of time at the particle physics level,
no activation function $\mathcal{E}(a)$, and no IAM sector split.
The other three forces maintain virial equilibrium in the stable
structures that decoherence produces. The weak force produces the
irreversibility that makes decoherence possible in the first place.

Three forces maintain stable equilibrium structures; one force
drives irreversible transitions between them. IAM's Law governs
the boundary between potential and actual; the weak force is the
crossing. The fact that exactly three of four forces satisfy the
equilibrium condition, and the one exception is the force whose
function requires it to not satisfy it, is structurally consistent
with IAM's Law identifying the boundary between quantum potential
and classical actuality as a thermodynamically costly one. Whether
this consistency constitutes independent evidence for the law or
merely a post-hoc structural observation is a question the authors
leave to the reader.

\subsection{Where IAM's Law Sits in the Landscape of Physics}

Every fundamental law in physics describes one side of a boundary.
Classical mechanics describes the post-classical regime --- the
deterministic evolution of definite states. General relativity
describes the geometry of actual spacetime --- the curved manifold
that results from the distribution of classical matter and energy.
Quantum mechanics describes the pre-classical regime --- the
evolution of superpositions before any definite outcome is selected.
Thermodynamics describes the statistical consequences of irreversible
processes in the classical regime. The Standard Model describes the
fields and particles that constitute classical matter once it exists.

Each of these frameworks is extraordinarily successful within its
tested domain. Each has a gap at the same location: the boundary
between quantum potential and classical actuality. Classical mechanics does
not explain why the macroscopic world is classical. General
relativity does not explain what it costs to produce the classical
matter whose energy curves spacetime. Quantum mechanics describes
superpositions but does not provide a mechanism for why superpositions
become definite outcomes. Thermodynamics states that entropy increases
but does not derive this from the dynamics of the quantum-to-classical
transition. The Standard Model describes particles but not the process
by which quantum fields become particles.

IAM's Law is a statement about the boundary itself.

The boundary is the decoherence event: the irreversible transition
from quantum superposition to definite classical outcome described
by Eqs.~\eqref{eq:decoherence}--\eqref{eq:rho_diagonal}. IAM's Law
states that this transition carries a thermodynamic cost
(Eq.~\eqref{eq:iams_law}), that this cost is irreversible, and that
it is encoded on the nearest available encoding surface. This single
statement has the following consequences, each connecting to an open
problem in an established framework:

\textbf{Connection to classical mechanics:} The macroscopic world
is classical because the actualization cost has been paid. Newton's
laws operate in the post-actualization regime. IAM's Law is the
mechanism that produces the regime in which Newton's laws hold.

\textbf{Connection to general relativity:} Jacobson's
1995 derivation~\citep{Jacobson1995} showed that the Einstein
equations follow from $\delta Q = T\,dS$ applied to local causal
horizons. The entropy functional was incomplete. IAM's Law identifies
the missing term: the accumulated Landauer cost of every irreversible
quantum-to-classical transition in the bulk. Adding it closes
Jacobson's derivation. The cosmological model derived here is the
observable consequence.

Beyond completing the entropy functional, IAM's Law also derives
both of Jacobson's remaining assumptions. The proportionality
$S \propto A$ follows from the causal accessibility of the horizon:
information produced by an irreversible decoherence event must be
encoded at the nearest 2D causal boundary because volume encoding
is causally inaccessible to external observers. The Bekenstein--Hawking
coefficient $\eta = 1/(4\ell_P^2)$ follows from the ratio
$1/4 = 2\pi/8\pi$, where $2\pi$ is the Euclidean Rindler cone angle
(the unique topological period eliminating the conical singularity at
the horizon) and $8\pi$ is the Einstein equation normalization (from
Gauss's law in three dimensions, doubled by the Bianchi identity
trace structure). Both factors are geometric facts about flat
spacetime independent of $S \propto A$. The same ratio determines
the minimum horizon area per decoherence event to be $4\ell_P^2$,
consistent with exactly one bit per $4\ell_P^2$ --- a self-consistency
that holds without observational input. These results are derived
in full in a companion paper~\citep{IAM_archive}. Jacobson's machine
--- \textit{specify an entropy functional, derive a gravity theory}
--- now fully specifies its own entropy functional from the same
causal geometry.

\textbf{Connection to thermodynamics:} The second law states that
entropy increases but does not derive this from the dynamics of the
quantum-to-classical transition. IAM's Law provides the mechanism:
entropy increases because every crossing from potential to actual
pays an irreversible Landauer cost to the nearest encoding surface.
The second law is a corollary. The activation function
$\mathcal{E}(a) = \exp(1 - 1/a)$, monotonically increasing and
bounded above, is the quantitative record of this irreversibility
applied across cosmic time.

\textbf{Connection to quantum mechanics:} The measurement problem
--- how quantum superpositions become definite classical outcomes ---
is reframed by IAM's Law from a foundational mystery into a
thermodynamic transaction. The transition occurs when a quantum
system interacts irreversibly with its environment, producing
classical information at Landauer cost $k_B T_H \ln 2$ per bit.
The transition is not instantaneous collapse but the irreversible
orthogonalization of environmental states described by
Eq.~\eqref{eq:decoherence}. IAM's Law does not solve the measurement
problem at the level of a UV-complete quantum gravity theory; it
identifies the thermodynamic cost of the transition and shows that
this cost, accumulated over cosmic history, has observable
cosmological consequences consistent with current data at $0.2\sigma$
across 17 independent MCMC chains.

\textbf{Connection to the cosmological constant:} The vacuum
baseline of information production --- quantum vacuum fluctuations
producing a constant, epoch-independent rate of classical information
--- maps directly to $\Lambda$ through the Jacobson mechanism. The
cosmological constant is the encoding cost of the vacuum baseline.
The structural complexity term $\beta_m\,\mathcal{E}(a)$ is the
encoding cost of gravitational structure formation above that
baseline. IAM's Law places both terms in the same entropy budget,
unified by the same cost per bit. The quantitative consequence is
derived in Section~\ref{sec:cc_baryon}: applying the Gibbons--Hawking
temperature ratio $\sqrt{\Omega_\Lambda} = 0.827$ closes the
cosmological constant to within $0.07\%$ of the observed value with
zero free parameters, and simultaneously selects the baryon asymmetry
$\eta = (6.115 \pm 0.037) \times 10^{-10}$ from the CMB alone.

The common thread is that each of these connections arises from the
same physical fact: every major framework in physics has a gap at
the quantum-to-classical boundary, and IAM's Law is a statement
about that boundary. This is not a claim that IAM's Law solves all
open problems in physics. It is a claim that those open problems
share a common location, that IAM's Law operates at that location,
and that its observable consequences are consistent with data to the
precision currently available.

The thermodynamic level of the derivation is complete. The
proportionality $S \propto A$ is derived from causal accessibility.
The coefficient $\eta = 1/(4\ell_P^2)$ is derived from the
geometric ratio $2\pi/8\pi$. The minimum encoding area $4\ell_P^2$
is derived from the same ratio. Jacobson's two assumptions reduce
to one derived result. The single remaining open question is why a
minimum resolvable length $\ell_P = \sqrt{\hbar G/c^3}$ exists at
all --- why the three constants $\hbar$, $G$, $c$ combine to set a
minimum Planck-scale area. That is a quantum gravity microstructure
question that thermodynamic approaches cannot address from within
their own framework. It is the same gap that exists in all
thermodynamic approaches to gravity, including Jacobson's original
derivation. IAM's Law is presented as a candidate law at this same
level of closure: derivationally self-consistent, with explicit
assumptions, and with the thermodynamic level fully closed. The
UV-complete microscopic account of Planck-scale discreteness
remains the open frontier.

\subsection{A Law, Not a Framework}

The major theories of physics are frameworks: collections of mutually
consistent postulates, equations, and structures that together
constitute a theory. General relativity requires the equivalence
principle, the metric tensor formalism, the geodesic equation, the
Einstein field equation, and the Bianchi identities --- assembled by
Einstein into a consistent structure, but not all consequences of a
single statement. The Standard Model requires a gauge group, a Higgs
mechanism, a fermion content, and 19 independently measured
parameters. Thermodynamics comprises four independently stated laws.
Quantum mechanics requires the Schr\"{o}dinger equation, the Born
rule, the measurement postulate, and the tensor product structure for
composite systems --- multiple independent inputs.

IAM is structurally different. Every result in this paper is a
consequence of one equation:
\begin{equation}
\Delta E_\mathrm{act} = k_B T_H \ln 2 \quad \text{per bit,}
\tag{\ref{eq:iams_law}}
\end{equation}
applied consistently at two scales using only established physics.
No additional postulates are introduced. The derivation chain is:

\begin{itemize}
\item The thermodynamic completion of the virial theorem
  (Section~\ref{sec:local}) follows from IAM's Law applied at each
  virialization event, with the virial temperature canceling exactly.

\item $\beta_m = \Omega_m/2$ (Eq.~\ref{eq:beta_derived}) follows
  from the virial $1/2$ partition as the balance sheet of IAM's Law.

\item $\mathcal{E}(a) = \exp(1-1/a)$ (Eq.~\ref{eq:Ea}) follows
  from integrating IAM's Law over cosmic time through the
  Cai--Kim first law.

\item $\mu(a) < 1$ (Eq.~\ref{eq:mu_exact}) follows from adding
  $S_\mathrm{info}$ to the Jacobson--Cai-Kim entropy functional.

\item $\Sigma(a) = 1$ (Eq.~\ref{eq:sigma_exact}) follows from
  photons propagating on null geodesics and paying no actualization
  cost.

\item The arrow of time follows from the cost being irreversible:
  $\mathcal{E}(a)$ is monotonically increasing because each payment
  is permanent.

\item Black holes as preferred early-universe encoding surfaces
  (Section~\ref{sec:bh}) follow from IAM's Law saturating the
  Bekenstein-Hawking bound locally before the cosmic horizon
  grows large enough to encode efficiently.

\item The coincidence problem resolution (Section~\ref{sec:arrow})
  follows from dark energy being driven by the same actualization
  cost as structure formation --- they are the same process at
  the same epoch by construction.

\item The sector-dependent Hubble split follows from the dual-sector
  separation that the cost produces between timelike and null
  worldlines.
\end{itemize}

Not one of these results required an additional free parameter or an
additional physical assumption beyond IAM's Law plus established
physics. This structural property --- that one equation generates its
consequences rather than requiring them to be assembled --- is what
the authors mean by calling IAM's Law a law rather than a framework.
The distinction is not a claim about final status; it is a
description of the derivational architecture.

Every boxed result in this paper traces back to
Eq.~\eqref{eq:iams_law} through a chain of established physics.
To the authors' knowledge, no additional free parameter or physical
assumption beyond Assumption~2 (Section~\ref{sec:assumptions}) is
required. The derivations are presented in sufficient detail for
this claim to be verified independently.

\subsection{Zero Free Parameters}

The entire IAM modification to $\Lambda$CDM is determined by
established physics and the independently measured matter density:
\begin{itemize}
\item $\mathcal{E}(a) = \exp(1-1/a)$: derived from horizon
  thermodynamics.
\item The $1/a$ exponent: derived from information surface density,
  confirmed to $1\%$ by the Sheth-Tormen mass function.
\item $\mathcal{E}(1) = 1$: set by definition.
\item $\beta_m = \Omega_m/2$: derived from the virial theorem.
\end{itemize}
IAM is a zero-parameter extension of $\Lambda$CDM. The framework
introduces no new fields, extra dimensions, or modifications to the
gravitational action beyond the identification of decoherence-produced
information as a contribution to horizon entropy.

IAM's Law has been applied to derive additional results across a range
of physical domains, including the M-$\sigma$ relation for
supermassive black holes~\citep{IAM_archive}, the lensing-dynamics
discrepancy in galaxy clusters, the phantom crossing signature in DESI
$w_0w_a$ fits as an artifact of fitting a dual-sector expansion history
with a single-sector model, the cosmological constant to within $0.07\%$
with zero free parameters (Section~\ref{sec:cc_baryon}), and the baryon
asymmetry $\eta = (6.115 \pm 0.037) \times 10^{-10}$ returned by the Planck
CMB likelihood with the BBN prior removed --- consistent with the
observed $6.137 \times 10^{-10}$ at $0.36\%$ with no nuclear physics
input (Section~\ref{sec:cc_baryon}). These results, with derivations
and timestamps, are archived at
\href{https://github.com/hmahaffeyges/IAM-Validation}{github.com/hmahaffeyges/IAM-Validation}
(Zenodo DOI:
\href{https://doi.org/10.5281/zenodo.18702042}{10.5281/zenodo.18702042}).

\subsection{The Mahaffey Number: IAM's Law in Dimensionless Form}
\label{sec:mahaffey_number}

IAM's Law states that every irreversible quantum-to-classical transition
pays $k_B T_H \ln 2$ per bit to the nearest encoding surface. Dividing
both sides of this statement by the Landauer floor yields a dimensionless
ratio that characterizes the operating regime of any physical system
relative to its encoding surface saturation point:

\begin{equation}
\mathcal{M} = \frac{E_{\rm drive}}{k_B T_{\rm local} \ln 2}
\label{eq:mahaffey}
\end{equation}

where $E_{\rm drive}$ is the energy source driving irreversible information
events in the system and $T_{\rm local}$ is the temperature of the local
encoding surface. $\mathcal{M}$ is the \textbf{Mahaffey Number} ---
a dimensionless figure characterizing how far above the Landauer floor
any informational actualization platform operates, following the
convention established by Reynolds (fluid flow regimes), Mach
(compressible flow), Froude (free surface flow), and Prandtl (heat
transfer) for dimensionless numbers that characterize physical operating
regimes. When $\mathcal{M} \geq \mathcal{M}_{\rm critical}$, the local
encoding surface saturates and a phase transition occurs.

$\mathcal{M}$ is not a new law. It is IAM's Law
(Eq.~\ref{eq:iams_law}) expressed in dimensionless form --- the natural
parameter that emerges when the drive energy is normalized by the Landauer
floor of the local encoding surface. Every application of IAM's Law
is a computation of $\mathcal{M}$ in domain-specific units:

\begin{itemize}
\item \textbf{Thermal systems} (semiconductor switching, superconducting
  qubit decoherence): $E_{\rm drive} = k_B T$, giving
  $\mathcal{M} = T/(T_{\rm local}\ln 2)$. Temperature is the lever.
  The performance scaling $A(T) \propto T^n$ is the observable consequence
  of $\mathcal{M}$ approaching the architecture-specific $\mathcal{M}_{\rm critical}$.

\item \textbf{Chemical systems} (biological cells): $E_{\rm drive} = \Delta G_{\rm ATP}$,
  $T_{\rm local} = T_{\rm body} = 310.15$~K (fixed). $\mathcal{M} =
  \Delta G_{\rm ATP}/(R_{\rm gas} T_{\rm body} \ln 2) = n_{\rm bio}/\ln 2 \approx 30.2$.
  Biology decoupled $E_{\rm drive}$ from $k_B T_{\rm local}$ by using
  chemical rather than thermal drive --- the evolutionary discovery that
  permits high-fidelity information maintenance at a fixed temperature far
  above the thermal floor.

\item \textbf{Cosmological} (gravitational decoherence): $E_{\rm drive}$
  is gravitational binding energy; $T_{\rm local} = T_H = \hbar H/(2\pi k_B)$.
  $\mathcal{M}$ integrated self-consistently over cosmic time with
  gravitational feedback \emph{is} the activation function
  $\mathcal{E}(a) = \exp(1 - 1/a)$ (Section~\ref{sec:activation}).
  The activation function is $\mathcal{M}$ solved with feedback.

\item \textbf{Gravitational saturation} (black hole formation):
  $E_{\rm drive} = Mc^2$; $T_{\rm local} = T_{\rm BH} = \hbar c^3/(8\pi G M k_B)$.
  $\mathcal{M} \propto (M/M_{\rm Planck})^2/\ln 2$, which is the
  Bekenstein--Hawking entropy in bits. The saturation condition
  $\mathcal{M} \geq 1/(4\ell_P^2)$ is black hole formation --- gravity
  is simultaneously the drive energy and the output of saturation.
\end{itemize}

The critical value $\mathcal{M}_{\rm critical}$ is architecture-specific,
set by the geometry of the local encoding surface: the same topological
period divided by geometric normalization structure that gives
$\eta = 2\pi/8\pi/\ell_P^2 = 1/(4\ell_P^2)$ for the gravitational case.
For superconducting qubits the critical condition is $T_1 \leq T_1^* = 0.65\,T_1^{\rm free}$
(the Substrate Inversion); for semiconductors it is $A_{\rm the semiconductor report} \leq A_{\rm floor}$
(the Dennard wall); for biological cells it is $\mathcal{A} = H(\beta)/H_{\rm min}(c) \geq 1.0$
(the $H_{\rm min}$ floor breach); for gravitational collapse it is
the Schwarzschild condition.

The universal statement is: IAM's Law leaves a fingerprint across every
physical domain. The fingerprint is $\mathcal{M}$. From the cosmic
horizon to the transistor junction to the human cell nucleus, every
irreversible information event is a computation of the same ratio ---
drive energy against Landauer floor, normalized by the architecture of
the local encoding surface.

% ─────────────────────────────────────────────────────────────────────────────
\section*{Acknowledgments}
% ─────────────────────────────────────────────────────────────────────────────

The thermodynamic framework of T.\ Jacobson~\citep{Jacobson1995} is
the foundation upon which this work is built. His 1995 derivation
identified gravity as thermodynamic and established the machine ---
specify an entropy functional, derive a gravity theory --- that IAM's
Law completes. The decoherence
framework of W.H.\ Zurek~\citep{Zurek2003} provided the local account
of the quantum-to-classical transition that the bridge extends to
cosmological scales. The author thanks R.\ Clausius~\citep{Clausius1870}
and R.\ Landauer~\citep{Landauer1961} for the tools that make this
connection possible.

% ─────────────────────────────────────────────────────────────────────────────
\section*{Data Availability}
% ─────────────────────────────────────────────────────────────────────────────

All MCMC chains, modified Fortran source, analysis scripts, and
prediction figures are publicly archived at
\url{https://github.com/hmahaffeyges/IAM-Validation} with canonical
Zenodo DOI
\href{https://doi.org/10.5281/zenodo.18702042}{10.5281/zenodo.18702042}.
Repository timestamps confirm all predictions were established prior
to any observational comparison.

\bibliographystyle{plainnat}
\bibliography{iam_law_v2}

\end{document}
